\section{\uppercase{Experimental setup}}\label{sec:experiments}

\subsection{Datasets}

\noindent
The evaluation uses five datasets, one simulated and four real, each in a distinct
role; Table~\ref{tab:datasets} lists their settings and splits.

The Webots TIAGO++ simulation supplies co-recorded WiFi and IMU trajectories. Its WiFi field is synthesized rather than measured, so the accuracy
below one meter it permits is optimistic relative to a real deployment. MSILN
site1/B1~\cite{microsoft2021indoorloc} is the real cross-session test: its
training, validation and test paths are recorded on separate days, so the test
fingerprints belong to sessions the model never observed. IMUWiFine
floor~4~\cite{nurpeiissov2022imuwifine} holds out paths within a single session,
testing within-session behavior on real WiFi--IMU data.
UJIIndoorLoc~\cite{torressospedra2014ujiindoorloc} is a WiFi fingerprinting
benchmark of unordered scans, testing the WiFi encoder in isolation. RoNIN
canonical~\cite{yan2019ronin} holds out test sequences from subjects absent from
training, testing the IMU encoder in isolation.

\begin{table*}[tb]
\caption{The five datasets. The Unit column names what each split count counts,
  as the datasets are not structured alike; \#APs is the number of WiFi access
  points where the benchmark fixes it. MSILN is a Microsoft Indoor Localisation
  collection; Webots is a TIAGO++ simulation.}\label{tab:datasets}
\centering
\begin{tabular}{|l|l|l|c|c|c|c|c|}\hline
Dataset & Modalities & Setting & \#APs & Unit & Train & Val & Test \\ \hline
Webots sim & WiFi, IMU & Simulation (lab) & 117 & paths & 11 & 3 & 3 \\ \hline
MSILN site1/B1 & WiFi, IMU & Real (cross-session) & 1419 & paths & 94 & 34 & 5 \\ \hline
IMUWiFine fl.4 & WiFi, IMU & Real (held-out paths) & 343 & paths & 40 & 20 & 20 \\ \hline
UJIIndoorLoc & WiFi & Real (isolation) & 520 & scans & 19937 & 1111 & n/a \\ \hline
RoNIN canonical & IMU & Real (isolation) & n/a & sequences & 73 & n/a & 32 \\ \hline
\end{tabular}
\end{table*}

\subsection{Reference Methods}

\noindent
One reference method is reported per claim, and each table column names it.

\begin{itemize}
\item \textbf{wlan\_localization}~\cite{wlanloc}: $k$-nearest-neighbor
  ($k$-NN) RSSI fingerprinting ($k=3$, distance-weighted); the WiFi reference for
  both the encoder-only and the end-to-end fusion comparison.
\item \textbf{ResNet1D}~\cite{yan2019ronin}: a $4.6$~M-parameter learned regressor
  over inertial windows; the IMU reference for the encoder-only inertial
  comparison.
\item \textbf{PDR-from-start}: IMU-only pedestrian dead reckoning with step
  detection, anchored at the first ground-truth waypoint; the classical inertial
  reference for the fusion comparison.
\item \textbf{IMUWiFine}~\cite{nurpeiissov2022imuwifine}: a four-layer LSTM
  that fuses WiFi and IMU; the learned-fusion reference for the comparison across
  sessions.
\end{itemize}

\subsection{Metrics}

\noindent
Two metrics are reported, both in meters. Let $\hat{p}_i \in \mathbb{R}^2$ be the
predicted planar position at sample $i$, $p_i \in \mathbb{R}^2$ the ground truth,
and $N_s$ the number of samples. Mean absolute error (MAE) is the mean Euclidean
position error,
%
\begin{equation}
  \mathrm{MAE} = \frac{1}{N_s} \sum_{i=1}^{N_s} \lVert \hat{p}_i - p_i \rVert_2 ,
  \label{eq:mae}
\end{equation}
%
and absolute trajectory error (ATE) is the root-mean-square Euclidean position
error, reported without alignment as the raw ATE,
%
\begin{equation}
  \mathrm{ATE} = \sqrt{ \frac{1}{N_s} \sum_{i=1}^{N_s}
    \lVert \hat{p}_i - p_i \rVert_2^2 } .
  \label{eq:ate}
\end{equation}
%
Umeyama-aligned ATE~\cite{umeyama1991}, which removes a rigid-body offset by
similarity alignment before measuring error, appears only in the limitations
discussion as contextual framing for the inertial encoder.

\subsection{Implementation and Training}

\noindent
The set transformer of Section~\ref{sec:method} is trained for $40$ epochs
(AdamW, learning rate $1.3\times10^{-3}$, weight decay $10^{-4}$; OneCycleLR;
Huber loss with $\delta=0.5$; gradient clipping at $1.0$; batch size $128$; $K=4$), with
a fixed seed of $42$, in PyTorch on a single NVIDIA Quadro P4000.
